\section{Integral Dependence and Valuations (In Progress)}

\subsection{Exercise - Complete}\label{5.1}

\begin{proof}
Note that by Exercise \ref{1.21}, since $f$ is surjective onto its image, $\Spec(A)$ is homeomorphic to $\Spec(f(A))$.
Since $f$ is an integral homomorphism, $B$ is an $f(A)$-algebra and is integral over $f(A)$. To see that $f^* : \Spec(A) \to \Spec(B)$ 
is a closed mapping, it suffices then to consider $f$ as an inclusion $A \xhookrightarrow{} B$.

A closed set of $\Spec(B)$ has the form $V(\mathfrak{b})$ for some ideal $\mathfrak{b}$ of $B$. By Exercise \ref{1.21}(iii),
$\overline{f^*(V(\mathfrak{b}))} = V(\mathfrak{b}^c)$ so that $f^*(V(\mathfrak{b})) \subset V(\mathfrak{b}^c)$. Note also that
$\mathfrak{b}^c = \mathfrak{b} \cap A$ since $A \subset B$. By Proposition 5.6, since $B$ is integral over $A$, $B/\mathfrak{b}$
is integral over $A/\mathfrak{b}^c$. Now, suppose that $\mathfrak{p} \in V(\mathfrak{b}^c)$. Then, $\overline{\mathfrak{p}} \in \Spec(A/\mathfrak{b}^c)$,
so by Theorem 5.10, there is some $\overline{\mathfrak{q}} \in \Spec(B/\mathfrak{b})$ such that $\overline{\mathfrak{q}} \cap A/\mathfrak{b}^c = \overline{\mathfrak{p}}$. 
In particular, $\mathfrak{q} \supset \mathfrak{b}$ and $\mathfrak{q} \cap A = \mathfrak{p}$. Thus, $\mathfrak{p} = f^{-1}(\mathfrak{q}) \in f^*(V(\mathfrak{b}))$. 
Hence, $f^*(V(\mathfrak{b})) = V(\mathfrak{b}^c)$ so that $f^*$ is a closed mapping.
\end{proof}

\subsection{Exercise - Complete}\label{5.2}

\begin{proof}
Since $\Omega$ is an integral domain, $f(A) \cong A/\ker(f)$ is also an integral domain. Thus, $\mathfrak{p} := \ker(f)$ is a prime ideal. By Theorem 5.10,
there is a prime ideal $\mathfrak{q}$ of $B$ such that $\mathfrak{q} \cap A = \mathfrak{p}$. By Proposition 5.6, $B/\mathfrak{q}$ is integral over
$A/\mathfrak{p}$. In particular, it suffices to prove the result when $A \subset B$ are integral domains with $B$ integral over $A$ and $f : A \to \Omega$ an
injection.

Let $\Sigma$ be the set of pairs $(C,\sigma)$, where $A \subset C \subset B$ and $\sigma : C \to \Omega$ is an embedding such that $\sigma$ restricts to $f$
on $A$. We can put a partial ordering on $\Sigma$ by letting $(C,\sigma) \le (C',\sigma')$ if $C \subset C'$ and $\sigma'$ restricts to $\sigma$ on $C$. Then,
$\Sigma$ is non-empty since it contains $(A,f)$. If we have a chain $(C_i,\sigma_i)$ of $\Sigma$, then $C = \bigcup_i C_i \subset B$ and $\sigma$ is defined on
$C$ by letting $\sigma(c) = \sigma_i(c)$ when $c \in C_i$; $c \in C$ implies $c \in C_i$ for some $i$, and this is well-defined by the compatibility requirement on 
restriction. Note that $\sigma$ is injective; if $\sigma(a) = \sigma(b)$, then since we're in a chain, they are identical on some common $i$, implying $a=b$ by the injectivity 
of each $\sigma_i$. Thus, $(C,\sigma) \in \Sigma$, so by Zorn's lemma, $\Sigma$ has maximal elements. If we have a maximal element of the form $(B,\sigma)$, we are
done, so suppose we have a maximal element $(C,\sigma)$, where $C \neq B$.

Suppose to the contrary that there is some $b \in B-C$; by assumption, $b$ is integral over $A$ and hence over $C$. Let $b$ satisfy $p(x) = \sum_i c_i x^i \in C[x]$.
Write $(\sigma p)(x) = \sum_i \sigma(c_i)x^i \in \Omega[x]$. Then, the composition
$$C[x] \to \Omega[x] \to \Omega[x]/((\sigma p)(x)) \cong \Omega$$
has kernel $(p(x))$, so we have an injection $C[x]/(p(x)) \to \Omega$ restricting to $\sigma$ on $C$. But $C[x]/(p(x)) \cong C[b]$, and by assumption
$C \subset C[b] \subset B$, where the first inclusion is proper. Thus, the existence of the induced map $C[b] \to \Omega$ violates the maximality of $\sigma$.
\end{proof}

\subsection{Exercise - Complete}\label{5.3}

\begin{proof}
Consider simple elements of $B' \otimes_A C$, i.e. elements of the form $b' \otimes c$, where $b' \in B'$ and $c \in C$. Since $B'$ is integral over $f(B)$
by assumption, there is some $p(x) = x^n + f(b_{n-1}) x^{n-1} + \cdots + f(b_0)$ with $b_i \in B$ such that $p(b') = 0$. Thus,
\begin{align*}
0 &= p(b') \\
0 &= p(b') \otimes c^n \\
0 &= (b')^n \otimes c^n + f(b_{n-1}) (b')^{n-1} \otimes c^n + \cdots + f(b_0) \otimes c^n \\
0 &= (b' \otimes c)^n + (f(b_{n-1}) \otimes c^1) (b' \otimes c)^{n-1} + \cdots + f(b_0) \otimes c^n.
\end{align*}
Thus, $b' \otimes c$ is integral over $f(B) \otimes C$, so every simple element of $B' \otimes C$ is integral over $f(B) \otimes C$.

In general, an element $x \in B' \otimes C$ is a finite sum of simple elements of $B' \otimes C$. By Corollary 5.3, the elements of
$B' \otimes C$ integral over $f(B) \otimes C$ form a subring of $B' \otimes C$. Since $x$ is a finite sum of simple elements, each of which is 
integral over $f(B) \otimes C$, $x$ is itself integral over $f(B) \otimes C$. Thus, $f \otimes 1 : B \otimes_A C \to B' \otimes_A C$ is integral
whenever $f$ is integral.
\end{proof}

\subsection{Exercise - Complete}\label{5.4}

No, if $B$ is integral over a subring $A$, $\mathfrak{n}$ is a maximal ideal of $B$, and $\mathfrak{m} = \mathfrak{n} \cap A$ is maximal in $A$, $B_{\mathfrak{n}}$
is not necessarily integral over $A_{\mathfrak{m}}$.

Let $A = k[x^2 - 1]$, $B = k[x]$, where $k$ is a characteristic $0$ field. Then, $f(t) = t^2 - ((x^2 - 1) + 1)$ shows that $x \in k[x]$ is integral over $k[x^2 - 1]$. 
The integral elements form a subring, so since everything in $k[x]$ is finitely generated by $k$ and $x$ through ring operations, $k[x]$ is integral over $k[x^2-1]$.

Now, let $\mathfrak{n} = (x-1)$ in $B$. We have a $k$-algebra surjective homomorphism $f : k[x] \to k$ by $x \mapsto 1$, and $(x-1)$ is the kernel,
so $k[x]/(x-1) \cong k$, a field. Thus, $\mathfrak{n}$ is maximal. Note that $(x^2 - 1) \subset k[x^2 - 1] \cap (x-1)$, but this ideal is maximal by 
the same argument, so $\mathfrak{m} = (x^2 - 1) = (x-1) \cap k[x^2-1]$ is maximal.

Note that $x+1 \in B - (x-1)$, so $\frac{1}{x+1} \in B_{\mathfrak{n}}$. Suppose $\frac{1}{x+1}$ is integral over $A_{\mathfrak{m}}$. Then,
$$0 = \left(\frac{1}{x+1}\right)^n + a_{n-1}\left(\frac{1}{x+1}\right)^{n-1} + \cdots + a_1 \frac{1}{x+1} + a_0$$
for some $a_i \in A_{\mathfrak{m}}$. Then, multiplying by $(x+1)^n$, we obtain
$$0 = 1 + a_{n-1}(x+1)^1 + \cdots + a_1 (x+1)^{n-1} + a_0 (x+1)^n.$$
Since $a_i \in A_{\mathfrak{m}}$ for each $i$, we can write $a_i = g_i(x)/h_i(x)$, where $g_i(x) \in k[x^2-1]$ and $h_i(x) \in k[x^2 - 1] - (x^2 - 1)$. In particular,
$h_i(-1)$ is the constant term of $h_i(x)$ since $(-1)^2 - 1 = 0$. However, $h_i(x) \not\in (x^2 - 1)$, so the constant term cannot be zero. Thus, for all $i$, we have
a well-defined $k$-algebra homomorphism $A_{\mathfrak{m}} \to k$ that sends $x \mapsto -1$; since $h_i(-1) \neq 0$, $g_i(-1)/h_i(-1) \in k$. Under this $k$-algebra homomorphism,
we obtain
$$0 = 1 + \frac{g_{n-1}(-1)}{h_{n-1}(-1)}(-1+1)^1 + \cdots + \frac{g_{1}(-1)}{h_{1}(-1)} (-1+1)^{n-1} + \frac{g_{0}(-1)}{h_{0}(-1)} (-1+1)^n,$$
implying that $0 = 1$, a clear contradiction. Thus, $\frac{1}{x+1} \in B_{\mathfrak{n}}$ is not integral over $A_{\mathfrak{m}}$.

\subsection{Exercise - Complete}\label{5.5}

\begin{proof}[(i) Proof.]
Since $x \in A$ is a unit in $B$, there is some $b \in B$ such that $xb = 1$. $B$ is integral over $A$, so there is some polynomial of
minimal degree $n > 0$ such that 
$$b^n + a_{n-1}b^{n-1} + \cdots + a_1 b + a_0 = 0,$$
where $a_i \in A$. Assume to the contrary that $n > 1$. Multiplying by $x$, we obtain
$$(xb)b^{n-1} + a_{n-1} (xb) b^{n-2} + \cdots + a_1 (xb) + a_0 x = 0,$$
but since $xb = 1$,
$$b^{n-1} + a_{n-1} b^{n-2} + \cdots + a_2 b + (a_1 + a_0 x) = 0.$$
Since $x \in A$, $a_1 + a_0 x \in A$, so this is a monic polynomial of $b$ (since $n > 1$) with coefficients in $A$. Since we chose $n > 1$ 
minimal, we clearly have a contradiction, so $n = 1$. In particular, $b$ satisfies a monic polynomial of the form $f(t) = t - a = 0$,
where $a \in A$, so $b = a \in A$, as desired.
\end{proof}

\begin{proof}[(ii) Proof.]
Let $\mathfrak{R}_B$ and $\mathfrak{R}_A$ be the Jacobson radicals of $B$ and $A$ respectively. We wish to show that $A \cap \mathfrak{R}_B = \mathfrak{R}_A$.
We will make use of Proposition 1.9, that $x \in \mathfrak{R} \iff 1 - xy$ is a unit in $R$ for all $y \in R$.

If $x \in \mathfrak{R}_B \cap A$, then $x \in A$ and $1 - xy$ is a unit in $B$ for all $y \in B$. Since $A \subset B$, $1 - xy$ is a unit in $B$
for all $y \in A$. In this case, $1 - xy \in A$, so by (i), since $1 - xy$ is a unit in $B$, it is also a unit in $A$. Thus, $1 - xy$ is a unit
in $A$ for all $y \in A$, so $x \in \mathfrak{R}_A$.

Conversely, suppose $x \not\in \mathfrak{R}_B \cap A$. If $x \not\in A$, then clearly $x \not\in \mathfrak{R}_A$, as desired. So, assume $x \in A$
so that $x \not\in \mathfrak{R}_B$. As $\mathfrak{R}_B$ is the intersection of maximal ideals of $B$, there is some maximal ideal $\mathfrak{m}$ of
$B$ such that $x \not\in \mathfrak{m}$. Since $\mathfrak{m}$ is prime, $\mathfrak{m}^c$ is also a prime ideal, so there is some maximal ideal 
$\mathfrak{m}_A$ of $A$ such that $\mathfrak{m}_A \supset \mathfrak{m}_B^c$. By Theorem 5.11 (``Going-Up''), there is a prime ideal $\mathfrak{q}$ such that
$\mathfrak{q} \supset \mathfrak{m}$ and $\mathfrak{q} \cap A = \mathfrak{m}_A$. Since $\mathfrak{m}$ is maximal, $\mathfrak{q} = \mathfrak{m}$, so 
$\mathfrak{m}_A = \mathfrak{m}^c$. Thus, $\mathfrak{m}^c$ is maximal. Because $\mathfrak{m}^c = A \cap \mathfrak{m}$, $x \not\in \mathfrak{m}$
implies $x \not\in \mathfrak{m}^c$. Then, since $\mathfrak{m}^c$ is maximal, $x \not\in \mathfrak{R}_A$.
\end{proof}

\subsection{Exercise - Complete}\label{5.6}

\begin{proof}
Since each $B_i$ is an integral $A$-algebra, each $B_i$ is integral over $f_i(A)$. Clearly, $f = (f_1,\ldots,f_n)$ makes $B := \prod_{i=1}^n B_i$ into
an $A$-algebra. So, we must show that $B$ is integral over $(\prod_{i=1}^n f_i)(A) = f(A)$. Let $b \in B$, and let $b_i$ be the component in $B_i$ under the
canonical projection map. Since $b_i$ is integral over $A$, $b_i$ satisfies some monic polynomial
$$b_i^m + f_i(a_{m-1})b_i^{m-1} + \cdots + f_i(a_1)b_i + f_i(a_0).$$
We can then define a polynomial $g_i(x)$ on $f(A)[x]$ as
$$g_i(x) = x^m + f(a_{m-1})x^{m-1} + \cdots + f(a_1)x + f(a_0).$$
If $\pi_i : B \to B_i$ is the canonical projection map,
$$\pi_i(g_i(b)) = \pi_i(b^m + f(a_{m-1})b^{m-1} + \cdots + f(a_0)) = b_i^m + f_i(a_{m-1})b_i^{m-1} + \cdots + f_i(a_0) = 0.$$
Thus, for each $i$, we have some monic polynomial $g_i \in f(A)[x]$ such that $\pi_i(g_i(b)) = 0$. Let $g = g_1 \cdots g_n \in f(A)[x]$; note that $g$
is monic since each factor is monic. Then, for all $i$,
$$\pi_i(g(b)) = \pi_i(g_1(b)) \cdots \pi_i(g_i(b)) \cdots \pi_i(g_n(b)) = \pi_i(g_1(b)) \cdots 0 \cdots \pi_i(g_n(b)) = 0,$$
so $g(b) = 0$. Hence, $b$ is integral over $f(A)$, so $B$ is an integral $A$-algebra.
\end{proof}

\subsection{Exercise - Complete}\label{5.7}

\begin{proof}
Suppose that $b \in B-A$ is integral over $A$. Then, it satisfies a monic polynomial in $A[x]$, so let's choose a monic polynomial of minimal degree $n > 0$
such that
$$b^n + a_{n-1}b^{n-1} + \cdots + a_1 b + a_0 = 0.$$
Since $b \not\in A$ by assumption, $n > 1$ since if $b + a_0 = 0$, $b = -a_0 \in A$. Rearranging terms, we have
$$b(b^{n-1} + a_{n-1}b^{n-2} + \cdots + a_1) = - a_0 \in A$$
Since $b \in B - A$ and $B - A$ is multiplicatively closed, we must have $b^{n-1} + a_{n-1}b^{n-2} + \cdots + a_1 \in A$; otherwise, both factors are in
$B-A$, so $-a_0 \in B-A$, a contradiction. In particular, we have that $b^{n-1} + a_{n-1}b^{n-2} + \cdots + a_1 = a \in A$, so
$$b^{n-1} + a_{n-1}b^{n-2} + \cdots + a_2 b + (a_1 - a) = 0,$$
where $a_1 - a \in A$. This violates the minimality of $n$, so $b \in B-A$ cannot be integral. Thus, $A$ is integrally closed in $B$.
\end{proof}

\subsection{Exercise}\label{5.8}

\begin{proof}[(i) Proof.]
Let $f,g \in B[x]$ be monic polynomials such that $fg \in C[x]$. Since $C$ is the integral closure of $A$ in $B$, $C$ is integrally closed in $B$ 
(Corollary 5.5). Since $B$ is an integral domain, it is a subring of its field of fractions, $K$ (Exercise \ref{3.9}). Thus, $C \subset B \subset K$,
so $fg \in K[x]$. Let $L$ be the splitting field of $fg$ over $K$; since $K \subset L$, $fg \in L[x]$ and splits into linear factors over $L[x]$. Thus,
$f(x) = \prod_{i=1}^n(x - \xi_i) \in L[x]$ and $g(x) = \prod_{j=1}^m(x - \eta_j) \in L[x]$. Since $\{\xi_i,\eta_j\}_{i,j}$ are roots of the polynomial
$fg \in C[x]$, then $\{\xi_i, \eta_j\}_{i,j} \subset L$ are all integral over $C$.

Let $E$ be the integral closure of $C$ in $L$. Then, $\{\xi_i, \eta_j\}_{i,j} \subset E$. Since $E$ is a subring of $L$ (Corollary 5.3), the coefficients of
$f(x) = \prod_{i=1}^n(x - \xi_i)$ and $g(x) = \prod_{j=1}^m(x - \eta_j)$ are all in $E$. Thus, $f(x),g(x) \in B[x] \cap E[x]$ ($B$ and $E$ viewed as 
subrings of $L$). Since $C$ is integrally closed in $B$, $C = B \cap E$, so $B[x] \cap E[x] = C[x]$. As desired, $f(x),g(x) \in C[x]$.
\end{proof}

\begin{proof}[(ii) Proof.]

\end{proof}

\subsection{Exercise}\label{5.9}

\subsection{Exercise}\label{5.10}

\subsection{Exercise}\label{5.11}

\subsection{Exercise}\label{5.12}

\subsection{Exercise}\label{5.13}

\subsection{Exercise}\label{5.14}

\subsection{Exercise}\label{5.15}

\subsection{Exercise}\label{5.16}

\subsection{Exercise}\label{5.17}

\subsection{Exercise}\label{5.18}

\subsection{Exercise}\label{5.19}

\subsection{Exercise}\label{5.20}

\subsection{Exercise}\label{5.21}

\subsection{Exercise}\label{5.22}

\subsection{Exercise}\label{5.23}

\subsection{Exercise}\label{5.24}

\subsection{Exercise}\label{5.25}

\subsection{Exercise}\label{5.26}

\subsection{Exercise}\label{5.27}

\subsection{Exercise}\label{5.28}

\subsection{Exercise}\label{5.29}

\subsection{Exercise}\label{5.30}

\subsection{Exercise}\label{5.31}

\subsection{Exercise}\label{5.32}

\subsection{Exercise}\label{5.33}

\subsection{Exercise}\label{5.34}

\subsection{Exercise}\label{5.35}
